% Copyright (c) 2026 Eugene Joseph M. Ragasa.
% SPDX-License-Identifier: Apache-2.0
% Author: Eugene Joseph M. Ragasa
% Formal proof completed: 2026-10-19
\documentclass[11pt]{article}
\usepackage[T1]{fontenc}
\usepackage{amsmath,amssymb,booktabs,tabularx}
\usepackage[margin=1in]{geometry}
\title{Proof Dossier: Omitted Weight Is Antitone in Support}
\author{Eugene Joseph M. Ragasa}
\date{Formal proof completed: 2026-10-19}
\begin{document}\maketitle
\section{Identity and status}
\begin{tabularx}{\textwidth}{@{}lX@{}}\toprule
Statement ID & \texttt{BG-FOURIER-THM-008} \\
Status & proved; no admissions \\
Lean declaration & \texttt{ProjectKoios.BandGauge.Fourier.}\newline\texttt{omittedWeight\_anti} \\
Owning module & \texttt{ProjectKoios/BandGauge/Fourier/FiniteReconstruction.lean} \\
Axioms & \texttt{propext}, \texttt{Classical.choice}, \texttt{Quot.sound} \\
\bottomrule\end{tabularx}
\section{Statement}
For pointwise nonnegative mode weights, enlarging a finite support cannot increase the total weight over omitted modes.
\section{Argument}
Support inclusion reverses under finite complementation. The omitted set for the larger support is therefore a subset of the omitted set for the smaller support, and adding nonnegative terms is monotone.
\section{Boundary and provenance}
The theorem does not assert that the supplied weights are Frobenius powers or physical locality measures. Lean checked the exact statement; the dossier and proof were AI-assisted.
\end{document}
